\documentclass[a4paper,titlepage]{article}
\usepackage[OT4,plmath]{polski}
%\usepackage[backend=biber]{biblatex}
\usepackage[margin=1in]{geometry}
\usepackage[noend]{algpseudocode}
\usepackage[page]{appendix}
\usepackage[usenames,dvipsnames,svgnames,table]{xcolor}
\usepackage[utf8]{inputenc}
\usepackage{adjustbox}
\usepackage{algorithm}
\usepackage{amsfonts}
\usepackage{amsmath}
\usepackage{amssymb}
\usepackage{amsthm}
\usepackage{array}
\usepackage{csquotes}
\usepackage{enumitem}
\usepackage{graphicx}
\usepackage{hyperref}
\usepackage{indentfirst}
\usepackage{longtable}
\usepackage{multirow}
\usepackage{parskip}
\usepackage{pifont}
\usepackage{setspace}
\usepackage{verbatimbox}
\usepackage{wrapfig}

\linespread{1.4}

%\addbibresource{Dokumentacja wstepna.bib}

\renewcommand*\appendixpagename{Załącznik}
\renewcommand{\qedsymbol}{$\square$}
\renewcommand{\algorithmiccomment}[1]{\hfill\textcolor{black!65}{\textit{#1}}}
\let\emptyset\varnothing

\makeatother
\setlength{\parindent}{24pt}
\theoremstyle{break}
\newtheorem*{uwaga}{Uwaga}
\newtheorem{definicja}{Definicja}[section]
\newtheorem{ozn}{Oznaczenie}[section]

\newcommand{\cmark}{\textcolor{ForestGreen}{\ding{51}}}
\newcommand{\xmark}{\textcolor{Maroon}{\ding{55}}}

%------------------------------------------------------------------------------

\title{Dokumentacja projektu}
\author{
    \href{mailto:bezap@student.mini.pw.edu.pl}{Patryk~Bęza}\\[0.7em]
    \href{mailto:kozakm@student.mini.pw.edu.pl}{Marek~Kozak}\\[0.7em]
    \href{mailto:malasnickik@student.mini.pw.edu.pl}{Krzysztof~Małaśnicki}\\[0.7em]
}
\date{\today}

\begin{document}

\makeatletter
\renewcommand{\ALG@name}{Algorytm}
\begin{titlepage}
\newcommand{\HRule}{\rule{\linewidth}{0.5mm}}
\center

\includegraphics[width=2.0cm]{img/mini}\\[1.5cm]
\textsc{\LARGE Politechnika Warszawska}\\[0.3cm]
\textsc{\Large Wydział Matematyki i~Nauk Informacyjnych}\\[1.5cm]
\textsc{\large Komputerowe wspomaganie technik kryminalistycznych}\\[0.2cm]
\textsc{\small Rok akademicki 2015/2016}\\[1.5cm]

\HRule \\[1cm]
{ \huge \bfseries \@title}\\[0.7cm]
\HRule \\[1.75cm]

\begin{minipage}[t]{0.4\textwidth}
\begin{flushleft}\large
\textsc{Autorzy:}\\[3mm]
\@author
\end{flushleft}
\end{minipage}
\begin{minipage}[t]{0.4\textwidth}
\begin{flushright}\large
\textsc{Wykładowca:}\\[3mm]
\href{mailto:m.szezynska@ise.pw.edu.pl}{dr~inż.~Magdalena~Szeżyńska}\\[1cm]
\end{flushright}
\end{minipage}
\vfill
{\large \today}

\end{titlepage}

%\maketitle
\tableofcontents
\clearpage

%------------------------------------------------------------------------------

\begin{abstract}

\begin{figure}[t]
    \centering
    \hspace{3em}
    \includegraphics[width=0.45\textwidth]{img/fingerprint}
\end{figure}

Niniejszy dokument powstał w~ramach dokumentacji projektu zespołowego z~przedmiotu \emph{\href{http://staff.elka.pw.edu.pl/~mszezyns/CAF/index.html}{Komputerowe wspomaganie technik kryminalistycznych}} w semestrze letnim roku akademickiego~2015/2016 na \emph{\href{http://www.mini.pw.edu.pl/}{Wydziale MiNI}~\href{https://www.pw.edu.pl/}{Politechniki Warszawskiej}}. Ma on za zadanie udokumentować ideę działania aplikacji, sposób jej użycia, wykorzystane narzędzia i biblioteki, wyniki testów funkcjonalnych i wydajnościowych oraz wnioski.

\end{abstract}

\clearpage

%------------------------------------------------------------------------------

\section{Wstęp}
\label{wstep}
\noindent W tym rozdziale przedstawiono krótki opis idei działania tworzonej aplikacji.

%------------------------------------------------------------------------------

\subsection{Cel aplikacji}

W ramach projektu powstanie aplikacja przeznaczona dla systemu~\emph{\href{https://www.microsoft.com/windows/}{Microsoft Windows}}, która będzie miała na celu... \emph{TODO}.

%------------------------------------------------------------------------------

\subsection{Idea działania aplikacji}

Rozwinięcie celu aplikacji \emph{TODO}...

%------------------------------------------------------------------------------

\section{Algorytm}

Do rozpoznawania twarzy zostanie zastosowana \emph{metoda analizy głównych składowych}~(ang.~\emph{Principal Component Analysis} -- PCA). Poniżej przedstawiono krótki opis tej metody.

%------------------------------------------------------------------------------

\subsection{Analiza głównych składowych}

Zakładamy, że zdjęcia treningowych twarzy są monochromatyczne i znormalizowane, tzn. twarz jest wyśrodkowana na zdjęciu i wszystkie zdjęcia mają jednakowe wymiary. Dla uproszczenia analizy, załóżmy, że zdjęcia mają wymiary $N\times N$, gdzie $N$ to wysokość i szerokość zdjęcia w pikselach.

\begin{description}[leftmargin=13em,style=nextline]
	\item[$M$] Liczba zdjęć twarzy treningowych.
	\item[$\Gamma_1, \Gamma_2, \dotsc, \Gamma_M$] $M$ wektorów o długości $N\cdot N$, reprezentujących $M$ zdjęć twarzy treningowych.
	\item[$\Psi = \frac{1}{M}\cdot\sum\limits_{n=1}^M\Gamma_n$] Wektor o długości $N\cdot N$, reprezentujący średnią z powyższego zbioru $M$ twarzy treningowych (twarz średnia z twarzy treningowych).
	\item[$\Phi_i=\Gamma_i-\Psi$] Wektor, będący różnicą między $i$-tym wektorem $\Gamma_i$, reprezentujący zdjęcie $i$-tej twarzy i wektorem reprezentującym średnią $\Psi$ wszystkich $M$ twarzy treningowych.
	\item[$C=\frac{1}{M}\cdot\sum\limits_{n=1}^M\Phi_n\Phi_n^T=AA^T$] Macierz kowariancji o wymiarach $N^2\times N^2$, gdzie $A=[\Phi_1\Phi_2\dotsc\Phi_M]$. Macierz $A$ ma wymiary~$N^2\times M$.
	\item[$u_1,u_2,\dotsc,u_K$] $K$ najlepszych wektorów własnych, tj. wektorów odpowiadającym $K$ największym wartościom własnym macierzy $C=AA^T$. Każdy z tych wektorów ma długość $M$.
\end{description}

Każda twarz w zbiorze treningowym może być zapisana jako kombinacja liniowa najlepszych $K$ wektorów własnych:
$$\Gamma_i-\Psi=\sum\limits_{j=1}^K w_j u_j\quad$$
gdzie:
$$w_j=u_j^T\Phi_i$$

%------------------------------------------------------------------------------

\section{Aplikacja}

Do stworzenia aplikacji zostanie wykorzystane środowisko programistyczne... \emph{TODO}.

%------------------------------------------------------------------------------

\subsection{Dane wejściowe}

\emph{TODO}

%------------------------------------------------------------------------------

\subsection{Dane wyjściowe}

\emph{TODO}

%------------------------------------------------------------------------------

\subsection{Instrukcja obsługi}

\emph{TODO}

%------------------------------------------------------------------------------

\section{Testy}

Ten rozdział opisuje wyniki przeprowadzonych testów funkcjonalnych i~wydajnościowych.

%------------------------------------------------------------------------------

\subsection{Dane testowe}

\emph{TODO}

%------------------------------------------------------------------------------

\subsection{Wyniki testów}

\emph{TODO}

%------------------------------------------------------------------------------

\newpage

\begin{appendices}

\section{Podział prac}
\subsection{Podział pracy pisemnej}

Poniższa tabela przestawia wykaz rozdziałów opracowanych przez poszczególnych członków zespołu.

\begin{table}[H]
\center
\begin{tabular}{p{2.5cm}|p{3.3cm}|p{3.3cm}|p{3.3cm}}
& Patryk Bęza & Marek Kozak & Krzysztof Małaśnicki \\\hline\hline
\parbox{3cm}{\ \\Opracowane \\rozdziały} & -- & -- & --\\
\end{tabular}
\end{table}

\subsection{Podział implementacji}

Poniższa tabela przestawia podział prac implementacyjnych w ramach projektu.

\begin{table}[H]
\center
\begin{tabular}{p{2.5cm}|p{3.3cm}|p{3.3cm}|p{3.3cm}}
& Patryk Bęza & Marek Kozak & Krzysztof Małaśnicki \\\hline\hline
\multirow{3}{*}{\parbox{3cm}{\ \\Opracowane \\funkcjonalności}} & $\bullet$ Bla & $\bullet$ Bla & $\bullet$ Bla\\
& $\bullet$ A & $\bullet$ B & $\bullet$ C \\
& $\bullet$ D & $\bullet$ E & $\bullet$ F \\
\end{tabular}
\end{table}

\end{appendices}

%------------------------------------------------------------------------------

%\clearpage
%\printbibliography[title=Bibliografia]

\end{document}
